\subsection{Verwandte Arbeiten}
\label{subsec:related-work}
Dieser Abschnitt stellt vier Arbeiten vor, die den Canny-Algorithmus auf
GPU- bzw. GPU-nahen Architekturen parallelisieren, mit besonderem Fokus
darauf, welche Speicheroptimierungen die jeweiligen Autoren für die
Non-Maximum Suppression einsetzen und begründen.

Der in diesem Praktikum implementierte Algorithmus basiert auf dem von
Canny \cite{canny1986} in \textit{A Computational Approach to Edge
Detection} vorgeschlagenen mehrstufigen Verfahren zur Kantendetektion,
bestehend aus Gaussian Blur, Gradientenberechnung, Non-Maximum
Suppression und Hysteresis Thresholding, dem die eigene Pipeline
unmittelbar folgt.

Song et al. \cite{song2024} implementieren in \textit{A Parallel Canny
Edge Detection Algorithm Based on OpenCL Acceleration} eine
OpenCL-basierte Variante und begründen den Einsatz von Local Memory
ausschließlich für die Faltungsstufen, da hier pro Pixel $n^2$ Nachbarn
gelesen werden; für die Non-Maximum Suppression, die nur zwei Nachbarn
pro Pixel benötigt, fehlt im Pseudocode dagegen jeder Lade- oder
Kachelungsschritt.

Ogawa et al. \cite{ogawa2010} stellen in \textit{Efficient Canny Edge
Detection using a GPU} eine frühe CUDA-Implementierung vor, bei der die
Non-Maximum Suppression ohne Kachelung direkt auf dem Global Memory
arbeitet. Der eigentliche Schwerpunkt liegt auf einer stack-basierten
Traversierung der nachgelagerten Hysterese-Stufe, die im Gegensatz zu
Ansätzen mit fester Traversierungstiefe auch lange, über mehrere
Bildkacheln verlaufende Kantenketten vollständig erfasst.

Horvath et al. \cite{horvath2023} implementieren in \textit{Canny Edge
Detection on GPU using CUDA} die Canny-Pipeline mit einem eigenen Kernel
pro Stufe. Die Shared-Memory-Implementierung der Non-Maximum-Suppression
orientiert sich am Tiling-Ansatz des Sobel-Kernels, benötigt aufgrund
der geringeren Nachbarschaftsgröße jedoch nur einen minimalen Halo-Rand.
Mittels NSIGHT Compute identifizieren die Autoren zudem
\texttt{hypot}/\texttt{atan2} als Hauptbottleneck ihres Sobel-Kernels;
deren Ersatz durch konstante Zuweisungen reduzierte die Kernel-Laufzeit
um den Faktor 8.0.

Yadav und Gupta \cite{yadav2025} vergleichen in \textit{Optimizing CUDA
Kernels for High-Performance Canny Edge Detection} vier
Optimierungsstufen anhand einer Roofline-Analyse, die sich jedoch
ausschließlich auf den Gaussian-Filter-Kernel bezieht. Bemerkenswert
ist, dass die Autoren die Non-Maximum Suppression konzeptionell nicht
dem Shared-Memory-Ansatz, sondern der Warp-Level-Optimierung zuordnen,
ohne diesen Kernel jedoch gesondert zu profilieren.